%
\renewcommand{\theequation}{\theenumi}
\begin{enumerate}[label=\arabic*.,ref=\thesubsection.\theenumi]
\numberwithin{equation}{enumi}

\item	In Fig. \ref{fig:tri_med_centroid}, $BE$ and $CF$ are the medians. Show that $\triangle GEF \sim \triangle GBC$.
%
\begin{figure}[!ht]
	\begin{center}
		\resizebox{\columnwidth}{!}{%Code by GVV Sharma
%December 10, 2019
%released under GNU GPL
%Drawing the median

\begin{tikzpicture}
[scale=2,>=stealth,point/.style={draw,circle,fill = black,inner sep=0.5pt},]

%Triangle sides
\def\a{5}
\def\b{6}
\def\c{4}
 
%Coordinates of A
\def\p{2.25}
\def\q{{sqrt(\c^2-\p^2)}}

%Labeling points
\node (A) at (\p,\q)[point,label=above right:$A$] {};
\node (B) at (0, 0)[point,label=below left:$B$] {};
\node (C) at (\a, 0)[point,label=below right:$C$] {};

%Foot of median

\node (D) at ($(B)!0.5!(C)$)[] {};
\node (E) at ($(A)!0.5!(C)$)[point,label=right:$E$] {};
\node (F) at ($(B)!0.5!(A)$)[point,label=left:$F$] {};

%Drawing triangle ABC
\draw (A) -- node[] {} (B) -- node[below, yshift=-5mm] {$\textrm{a}$} (C) -- node[] {} (A);

%Drawing medians BE and CF
\draw (B) -- (E);
\draw (C) -- (F);

%Drawing EF
\draw (E) -- (F);

%Labeling sides
\node [right] at ($(A)!0.5!(E)$) {$\frac{b}{2}$};
\node [right] at ($(C)!0.5!(E)$) {$\frac{b}{2}$};
\node [left] at ($(B)!0.5!(F)$) {$\frac{c}{2}$};
\node [left] at ($(A)!0.5!(F)$) {$\frac{c}{2}$};

%Centroid
\node (G) at ($(B)!0.67!(E)$)[label={[shift={(0.8,-0.5)}]$G$}] {};

%Angles
\tkzMarkAngle[fill=green!50,size=.3](F,E,B)
\tkzMarkAngle[fill=green!50,size=.3](C,B,E)
%
\tkzMarkAngle[fill=red!50,size=.3](C,F,E)
\tkzMarkAngle[fill=red!50,size=.3](F,C,B)
%
\tkzMarkAngle[fill=blue!50,size=.3](E,G,F)
\tkzMarkAngle[fill=blue!50,size=.3](B,G,C)

\tkzLabelAngle[pos=0.4](F,G,E){$\theta$}
\tkzLabelAngle[pos=0.4](B,G,C){$\theta$}
\tkzLabelAngle[pos=-0.4](F,E,G){$\alpha$}
\tkzLabelAngle[pos=0.65](C,B,E){$\alpha$}
\tkzLabelAngle[pos=0.65](C,F,E){$\beta$}
\tkzLabelAngle[pos=0.65](F,C,B){$\beta$}

\end{tikzpicture}

}
	\end{center}
	\caption{Similar Triangles}
	\label{fig:tri_med_centroid}	
\end{figure}
\\
\solution $\because EF \parallel BC$, alternate interior angles are equal.  Hence, corresponding angles of both triangles are equal and the triangles are similar.
\item Show that

%	Using Fig. \ref{ch2_median_ratio_val}, 
	\begin{align}
\label{eq:tri_med_centroid_ratio}
	\frac{GB}{GE} = \frac{GC}{GF} = 2
	\end{align}
%
\solution $\because \triangle GEF \sim \triangle GBC$ and $\frac{EF}{BC} =2$, \eqref{eq:tri_med_centroid_ratio} follows.
%
\item In Fig. \ref{fig:tri_med_meet}, AG is extended to meet $BC$ at $D$.	Show that 
	\begin{align}
\label{eq:tri_med_centroid_ratio_AD}
\frac{GA}{GD} = 2
	\end{align}
\begin{figure}[!ht]
	\begin{center}
		\resizebox{\columnwidth}{!}{%Code by GVV Sharma
%December 10, 2019
%released under GNU GPL
%Drawing the median

\begin{tikzpicture}
[scale=2,>=stealth,point/.style={draw,circle,fill = black,inner sep=0.5pt},]

%Triangle sides
\def\a{5}
\def\b{6}
\def\c{4}
 
%Coordinates of A
\def\p{2.25}
\def\q{{sqrt(\c^2-\p^2)}}

%Labeling points
\node (A) at (\p,\q)[point,label=above right:$A$] {};
\node (B) at (0, 0)[point,label=below left:$B$] {};
\node (C) at (\a, 0)[point,label=below right:$C$] {};

%Foot of median

\node (D) at ($(B)!0.5!(C)$)[point,label=below:$D$] {};
\node (E) at ($(A)!0.5!(C)$)[point,label=right:$E$] {};
\node (F) at ($(B)!0.5!(A)$)[point,label=left:$F$] {};

%Drawing triangle ABC
\draw (A) -- node[] {} (B) -- node[below, yshift=-5mm] {$\textrm{a}$} (C) -- node[] {} (A);

%Drawing medians AD, BE and CF
\draw (B) -- (E);
\draw (C) -- (F);
\draw (A) -- (D);

%Drawing EF
\draw [dashed] (E) -- (F);

%Centroid
\node (G) at ($(B)!0.67!(E)$)[label={[shift={(0.8,-0.5)}]$G$}] {};

%Labeling sides
\node [right] at ($(A)!0.5!(E)$) {$\frac{b}{2}$};
\node [right] at ($(C)!0.5!(E)$) {$\frac{b}{2}$};
\node [left] at ($(B)!0.5!(F)$) {$\frac{c}{2}$};
\node [left] at ($(A)!0.5!(F)$) {$\frac{c}{2}$};
\node [below] at ($(E)!0.5!(G)$) {$p$};
\node [below] at ($(B)!0.5!(G)$) {$2p$};
\node [above right] at ($(F)!0.5!(E)$) {$P$};

%\node (G) at ($(B)!0.67!(E)$)[label={[shift={(-0.8,-0.5)}]$G_1$}] {};

%
\end{tikzpicture}

}
	\end{center}
	\caption{Similar Triangles}
	\label{fig:tri_med_meet}	
\end{figure}

\solution  From Theorem \ref{them:tri_med_bpt}, $ EF || BC$.  Hence,
%
\begin{align}
\triangle APE &\sim \triangle ADC \quad (AAA)
\\
\implies AP &= PD = \frac{AD}{2}
\\
\implies  AG - GP &= GP + GD 
\\
\text{or, }  GP=\frac{AG - GD}{2} 
\label{eq:tri_med_centroid_ratio_AD1}
\end{align}
%
Similarly, 
%
\begin{align}
\triangle PGE &\sim \triangle BGD  \quad (AAA)
\\
\implies \frac{GP}{GD} &= \frac{GE}{GB} = \frac{1}{2}
\\
\text{or, }  GP=\frac{GD}{2} 
\label{eq:tri_med_centroid_ratio_AD2}
\end{align}
using \eqref{eq:tri_med_centroid_ratio}. From \label{eq:tri_med_centroid_ratio_AD1}
and \label{eq:tri_med_centroid_ratio_AD2},
%
\begin{align}
GP=\frac{AG - GD}{2} &=\frac{GD}{2}  
\\
\implies  \frac{AG}{GD} &=2
\end{align}
%%
\item Show that 
%
\begin{align}
BD = DC
\end{align}
%%
\item Show that the medians of a triangle meet at a point.
\\
\solution In Fig. \ref{fig:tri_med_meet}, $BE$ and $CF$	are medians that meet at $\vec{G}$.  Hence, 
from \eqref{tri_med_centroid_ratio}, 
%
\begin{align}
\label{eq:tri_centroid_g}
\vec{G} = \frac{2\vec{B}+\vec{E}}{3}
\end{align}
%
If the median $AD$ meets $BE$ in $\vec{G}_1$, from Theorem \eqref{tri_med_centroid_ratio}, 
%
\begin{align}
\vec{G}_1 = \frac{2\vec{B}+\vec{E}}{3}
\label{eq:tri_centroid_g1}
\end{align}
%
From  \eqref{eq:tri_centroid_g}
 and \eqref{eq:tri_centroid_g1}, $\vec{G} = \vec{G}_1$.  Thus, the medians meet at $\vec{G}$.
\item The centroid divides the median in the ratio $2:1$.
\label{them:tri_med_ratio}



%\item
%	In Fig. \ref{ch2_median_3}, the line $AO$ cuts $EF$ at $G$ and is extended to meet the side $BC$ at $D$.  Show that 
%	%
%	\begin{equation}
%	\frac{OA}{OD} = 2.
%	\end{equation}
%	%  
%
%\begin{figure}[!ht]
%	\begin{center}
%		
%		%\includegraphics[width=\columnwidth]{./figs/ch2_median_3}
%		%\vspace*{-10cm}
%		\resizebox{\columnwidth}{!}{\begin{tikzpicture}
[scale=2,>=stealth,point/.style={draw,circle,fill = black,inner sep=0.5pt},]

\node (E) at (1.5, 1.5)[point,label=above :$E$] {};
\node (F) at (-1.5, 1.5)[point,label=above :$F$] {};
\node (A) at (0, 3)[point,label=above :$A$]{};
\node (B) at (-3, 0)[point,label=below left:$B$]{};
\node (C) at (3, 0)[point,label=below right:$C$]{};
\node (O) at (0,1)[point,label=below right :$O$] {};
\node (D) at (0,0)[point,label=below :$D$] {};

\draw (B)--(A);
\draw (A)--(C);
\draw (C)--(B);
\draw (B)--(E);
\draw (C)--(F);
\draw (E)--(F);
\draw (A)--(D);

\node [above] at (-1.7,1.7) {$c$};
\node [above] at (1.7,1.7) {$b$};
\node [above] at (-2.5,.75) {$c/2$};
\node [above] at (-1,2.1) {$c/2$};
\node [above] at (2.5,.75) {$b/2$};
\node [above] at (1,2.1) {$b/2$};
\node [below] at (1,1.3) {$p$};
\node [below] at (-1,1.3) {$q$};
\node [above] at (-.7,.5) {$2p$};
\node [above] at (.7,.5) {$2q$};
\node [above] at (0,-.5) {$a$};
\node [above] at (0,1.9) [point,label=above right :$r$] {};
\node (G) at (0,1.5)[point,label=above right :$G$] {};

\tkzMarkAngle[size=.3](F,O,B);
\tkzMarkAngle[size=.3](C,O,E);
\draw (-0.5,1) node{$\theta$};
\draw (0.5,1) node{$\theta$};

\end{tikzpicture}}
%	\end{center}
%	\caption{Similar Triangles and Median}
%	\label{ch2_median_3}	
%\end{figure}
%
%
%\item
%	In Fig. \ref{ch2_median_ratio_val}, show that
%	\begin{equation}
%	ar\brak{\triangle AFE} = 	\frac{1}{4} ar\brak{\triangle ABC}
%	\end{equation}
%	\label{ch2_small_triang}
%%
%\solution We have
%%%
%\begin{align}
%ar\brak{\triangle AFE} &= \frac{1}{2} \frac{b}{2}\frac{c}{2}\sin A = \frac{1}{4}.\frac{1}{2}bc \sin A \\
%&=\frac{1}{4}ar\brak{\triangle ABC}
%\end{align}
%%
%		%\includegraphics[width=\columnwidth]{./figs/ch2_median_ratio_val}
%		%\vspace*{-10cm}
%
%%
%\begin{figure}[!ht]
%	\begin{center}
%		
%		\resizebox{\columnwidth}{!}{\begin{tikzpicture}
[scale=2,>=stealth,point/.style={draw,circle,fill = black,inner sep=0.5pt},]

\node (E) at (1.5, 1.5)[point,label=above :$E$] {};
\node (F) at (-1.5, 1.5)[point,label=above :$F$] {};
\node (A) at (0, 3)[point,label=above :$A$]{};
\node (B) at (-3, 0)[point,label=below left:$B$]{};
\node (C) at (3, 0)[point,label=below right:$C$]{};
\node (a) at (0,0)[point,label=below :$a$] {};
\node (O) at (0,1)[point,label=below :$O$] {};

\draw (a)--(B);
\draw (B)--(A);
\draw (A)--(C);
\draw (C)--(a);
\draw (B)--(E);
\draw (C)--(F);
\draw (E)--(F);
\node [above] at (-1.7,1.7) {$c$};
\node [above] at (1.7,1.7) {$b$};
\node [above] at (-2.5,.75) {$c/2$};
\node [above] at (-1,2.1) {$c/2$};
\node [above] at (2.5,.75) {$b/2$};
\node [above] at (1,2.1) {$b/2$};
\node [below] at (1,1.3) {$p$};
\node [below] at (-1,1.3) {$q$};
\node [above] at (-.7,.5) {$k_{p}$};
\node [above] at (.7,.5) {$k_{q}$};

\tkzMarkAngle[size=.3](F,O,B);
\tkzMarkAngle[size=.3](C,O,E);
\draw (-0.5,1) node{$\theta$};
\draw (0.5,1) node{$\theta$};

\end{tikzpicture}}
%	\end{center}
%	\caption{$O$ divides medians in the ratio $2:1$}
%	\label{ch2_median_ratio_val}	
%\end{figure}
%%
%
%\solution Using the sine formula and \eqref{ch2_small_triang}, areas of some triangles in Fig. \ref{ch2_median_ratio_val} are listed in the following table
%\begin{center}
%\begin{tabular}{|c|c|}
%	\hline  Triangle& Area  \\ 
%	\hline  OFE & $\frac{1}{2}pq \sin \theta $ \\ 
%	\hline  BOF &  $\frac{k}{2}pq \sin \theta$\\ 
%	\hline  COE & $\frac{k}{2}pq \sin \theta$ \\ 
%	\hline  BOC&  $\frac{k^2}{2}pq \sin \theta$\\
%	\hline  BOC&  $\frac{1}{2} ar \brak{\triangle ABC}$\\  
%	\hline 
%\end{tabular} 
%\end{center}
%%
%where we have used the fact that 
%
%\begin{equation}
%\sin \angle BOC = \sin \brak{180^{\degree}-\theta} = \sin \theta
%\end{equation}
%%
%Since $BE$ is the median
%\begin{align}
%\begin{split}
%ar\brak{\triangle BEC} &= \frac{1}{2} ar\brak{\triangle ABC}  
%\\
%& = ar\brak{\triangle BOC} + ar\brak{\triangle COE}\\
%ar\brak{\triangle BEA} &= \frac{1}{2} ar\brak{\triangle ABC}
%\\
%& = ar\brak{\triangle AFE} + ar\brak{\triangle BOF} +  ar\brak{\triangle FOE}  
%\end{split}
%\end{align}
%%
%Substituting the respective values from the above table,
%\begin{align}
%\begin{split}
%\frac{1}{2} ar\brak{\triangle ABC} &= \frac{k^2}{2}pq \sin \theta + \frac{k}{2}pq \sin \theta
%\\
%\frac{1}{2} ar\brak{\triangle ABC} &= \frac{k}{2}pq \sin \theta + \frac{1}{2}pq \sin \theta + \frac{1}{4} ar\brak{\triangle ABC}
%\end{split}
%\end{align}
%%
%Simplifying the above,
%%
%\begin{align}
%k(k+1) = 2(k+1)
%\end{align}
%%
%Since $k \neq -1$, $k = 2$ and the proof is complete.
%

\end{enumerate}
